% =============================================================================
% Conclusion Section --- Film Production Incentives Gravity Model
%
% Slots into main document under a \section{Conclusion} header.
% This file is written as continuous prose with paragraph breaks rather
% than subsections, following standard applied-economics convention.
%
% Required packages: natbib (for \citep, \citet)
%
% Bibliography keys referenced:
%   borusyak2024, callaway2021
%
% Section labels referenced:
%   sec:limitations
% =============================================================================

This paper has examined whether national film production incentive schemes attract international film production, using a gravity model framework applied to two complementary datasets: IMDb filming-location data covering the top 600 grossing films per year between 1990 and 2023, and OECD--WTO Balanced Trade in Services data on bilateral audiovisual services exports from 2005 to 2023. Identification proceeds through a sequence of specifications, beginning with PPML estimation of the standard gravity equation and progressing to OLS with country-pair fixed effects, which absorbs all time-invariant bilateral heterogeneity and identifies incentive effects exclusively from within-pair temporal variation. An event study examines the dynamic structure of the response and tests for pre-existing trends.

Four main findings emerge. First, the gravity framework fits bilateral film production flows well but the determinants differ substantially from those documented for goods trade: distance elasticities are smaller, common language effects are larger, contiguity effects are positive but more modest, and colonial ties and regional trade agreements are negative for film production rather than positive as they are for goods. The framework fits the BaTIS audiovisual services data more closely than the IMDb data, with the BaTIS gravity estimates aligning more closely with conventional goods-trade benchmarks. Second, incentive schemes do attract international production, but the cross-sectional estimate of approximately 18\% more inward production in countries with active schemes attenuates to roughly 9\% once country-pair fixed effects absorb time-invariant bilateral heterogeneity. The within-pair effect remains economically meaningful and is robust across subsamples and across the IMDb and BaTIS datasets, but the gap between the two estimates suggests that approximately half of the cross-sectional correlation reflects selection: countries with established film industries were systematically more likely to adopt incentives. Third, the timing of adoption matters but the duration of operation does not. Early adopters (pre-2005) attract substantially more production than late adopters (approximately 26\% more films and 45\% more spending), while incentives show no strong evidence of a cumulative "cluster effect". Fourth, several results change materially when the five major Anglosphere countries are excluded from the sample, with the destination incentive effect weakening notably and some specifications losing significance entirely. This suggests that the destination effect is concentrated in the Anglo production network rather than generalising uniformly across all bilateral pairs.

Taken together, these findings sit uneasily with both ends of the existing debate over film incentive effectiveness. The within-pair estimate is large enough to indicate a genuine treatment effect rather than pure selection, contradicting the strongest sceptical reading of the literature. But the magnitude is substantially smaller than cross-sectional estimates would imply, and the early-adopter premium suggests that the headline returns reported by industry-commissioned evaluations reflect a less competitive policy landscape that no longer exists. Furthermore, the incentive effect appears to be a one-time level shift attributable to entering the international subsidy competition, rather than a compounding advantage that grows with time.

The main policy implication of this work is for decision makers which are  considering adopting incentive schemes today. A country considering adoption today does not face the policy environment that delivered the 26\% film-count premium to early adopters; it faces a substantially more saturated landscape in which the marginal effect of an additional scheme is plausibly smaller than even the late-adopter estimates reported here. The estimates in this paper should therefore be regarded as upper bounds on the marginal effect of new adoption. Additionally, none of this speaks to the seperate issue of whether incentive schemes generate enough additional economic activity to justify their cost.

Several limitations of the analysis warrant emphasis and are discussed in full in Section~\ref{sec:limitations}. The country-pair fixed effects design absorbs time-invariant bilateral heterogeneity but cannot rule out time-varying confounders. The two-way fixed effects estimator may also be biased under treatment effect heterogeneity across cohorts, a setting that the Model~7 finding on early versus late adopters confirms is empirically relevant. The IMDb sampling frame is biased toward Anglo-language theatrical releases and  under-represents the streaming-era production patterns that now dominate the industry, raising questions about how well the historical estimates generalise to contemporary production decisions.

Two extensions seem particularly worth pursuing. First, re-estimation using cohort-robust difference-in-differences estimators such as the imputation approach of \citet{borusyak2024} or the group-time average treatment effect estimator of \citet{callaway2021} would test whether the within-pair destination effect of approximately 9\% is sensitive to the specific weighting structure of two-way fixed effects under staggered adoption with heterogeneous treatment effects. Second, a more comprehensive treatment of the different incentive schemes available, both by including subnational incentive schemes in the analysis and by more finely catagorising the characteristics of different incentive schemes so that they can be appropriately compared to one another. The inclusion of subnational incentive schemes into the international analysis would address one of the most significant data limitations of the present design, namely that many of the most generous schemes operate at the subnational level. These extensions would help address the question of how effective incentive schemes actually are at attracting international production and would help identify which aspects of particular incentive schemes are most effective at doing so, providing a more complete picture for policymakers who are thinking about implementing such schemes.
